\begin{abstract}
Teachers can find thousands of openly licensed textbooks, but judging whether one is good, and whether it can actually be reused, still rests on opinion. A systematic review of 14 published OER quality rubrics found that five specify no rating scale, only three provide a detailed scoring guide, and only six report any empirical testing. This position paper argues that open educational resources need an objective score: one anchored by descriptors and evidence, reported as separate quality and reusability halves, reproducible by human and AI-agent reviewers, and validated before anyone acts on it. We propose the OpenBook QR Score, ten metrics worth ten points each, five for quality (Q) and five for reusability (R), reported in the form QR 76 (Q 37 / R 39), with a decision rule that stops a high total from hiding a failed metric. As an illustration, not a validation, AI-agent reviewers scored each of ten widely used open textbooks in two independent passes, using about 86,000 tokens and two minutes per title; totals ranged from 60 to 81, and the passes agreed within one point on 88 percent of metric scores. Metadata was the weakest metric (mean 5.5), editable source the strongest (8.9), and eight of ten titles carry a NonCommercial clause that the descriptors do not yet score consistently. Scripted checks on 101 MkDocs book projects apply part of a rubric in seconds per book. We set out the reliability, validity and gaming studies required before scores guide adoption.
\end{abstract}

\noindent\textbf{Keywords:} open educational resources, OER quality, rubric, reusability, AI agents, validation
